\chapter{Polygons: the PolyYOLO radial representation}
\label{ch:polygons}

The polygon head is the most distinctive part of this codebase, so it gets its own
chapter. The goal: predict an object's outline as 24 points, cheaply enough to be just
three small per-pixel heads.

\section{The radial idea}
Instead of storing 24 arbitrary $(x,y)$ points, PolyYOLO fixes the \emph{angles} and
predicts only the \emph{radii}. Imagine standing at the object's box center and shooting
24 rays outward, one every $15^\circ$ ($360^\circ / 24$). For each ray, predict how far
the boundary is. That distance is the only number you need per vertex --- the angle is
implied by the bin index.

\guidefig{fig_radial.png}{1.0}{The radial encoding. Left: 24 angular bins ($15^\circ$
each) around the box center; the green segment in each occupied bin is the encoded radius
to the polygon boundary (blue). Right: the per-bin target --- a radius (\code{dist}) and a
validity flag (\code{conf}) for each of the 24 bins. This is exactly the target built for
the loss.}

\begin{hood}
Each vertex is actually \emph{three} numbers, so the target is \code{[24 $\times$ 3] = 72}
values per object, interleaved as \code{[dist, angle, conf]}:
\begin{itemize}[leftmargin=1.4em,itemsep=1pt]
  \item \code{dist} --- the radius (how far the boundary is along this ray);
  \item \code{angle} --- a \textbf{sub-bin offset} in $[0,1)$: the exact vertex angle is
        $(i + \text{offset})\cdot15^\circ$, so the vertex isn't snapped to the bin center;
  \item \code{conf} --- 1 if this bin holds a vertex, 0 if empty.
\end{itemize}
This is built in \file{losses/tal\_loss.py}\,:\,\code{\_polygon\_loss}. The three model
heads \code{poly\_dist}, \code{poly\_angle}, \code{poly\_conf} predict these channels.
\end{hood}

\begin{teacher}
Why a sub-bin offset instead of just the bin index? With $15^\circ$ bins, snapping every
vertex to its bin center would quantize the outline coarsely. The offset recovers the
\emph{exact} angle within the bin, so the 24 vertices can sit anywhere on the boundary ---
much smoother outlines for the same 24 bins.
\end{teacher}

\section{Why confidence sees every bin}
Not every object has a vertex in every direction (a thin object leaves many bins empty).
The \code{conf} channel is the gate that decides, at decode time, which bins emit a vertex
(threshold $0.4$).

\begin{why}
The \code{conf} loss is averaged over \textbf{all 24 bins} (occupied $\rightarrow$ 1, empty
$\rightarrow$ 0), while \code{dist} and \code{angle} are averaged over occupied bins only.
The reason is subtle and was found the hard way: \code{conf} is the gate, so it must be
\emph{trained on negatives} (empty bins) to learn to output low confidence there. An
earlier version masked \code{conf} to occupied bins only --- so empty bins never got a
``you should be off'' signal, their confidence drifted above the $0.4$ threshold, and the
decoder emitted spurious vertices: the ``spiky star'' artifacts seen in validation
overlays. \code{dist}/\code{angle} stay masked because a radius/angle on an empty bin is
undefined --- there is no vertex to regress toward.
\end{why}

\section{Arc-length resampling: making bins track real shapes}
There is a catch. Datasets often store polygons with few vertices --- a rectangle might be
just its 4 corners. Drop 4 corners into 24 angular bins and only $\sim$4 bins get a
vertex; the radial target becomes a diamond, and the long edges (which cross many empty
bins) are never learned.

\guidefig{fig_resample.png}{0.92}{The fix, run through the real
\code{resample\_polygons}. A 4-corner rectangle (left) is resampled to 64 points spaced
uniformly \emph{along} the contour (right). The new points lie on the edges, so every
angular bin the boundary crosses now gets a vertex --- the rectangle encodes as a
rectangle, not a diamond.}

\begin{hood}
\code{resample\_polygons} (in \file{data\_pipeline/augmentations.py}) walks the closed
contour and places \code{K=64} points at equal arc-length intervals, interpolating
\emph{on the edges} --- it does not just subsample stored vertices. This runs at decode
time (\code{parser.resample\_points: 64}), so every downstream stage carries a compact
\code{[N, 128]} polygon instead of the raw stored width (which can be thousands of
columns).
\end{hood}

\section{Three formats across the stack}
The same polygon wears three different ``costumes'' as it moves through the code. Keeping
them straight is essential.

\begin{center}
\begin{tikzpicture}[node distance=6mm and 14mm,font=\small]
  \node[blkn,align=left] (a) {\textbf{1. Dataset / TFDS}\\\scriptsize flat $xy$ interleaved\\\scriptsize normalized, $-1$ padded\\\scriptsize \code{[N, 3972]}};
  \node[blka,right=of a,align=left] (b) {\textbf{2. Training target}\\\scriptsize radial \code{[dist,angle,conf]}$\times24$\\\scriptsize \code{[N, 72]}\\\scriptsize origin = box center};
  \node[blkg,right=of b,align=left] (c) {\textbf{3. Cartesian}\\\scriptsize pixel $(x,y)$ vertices\\\scriptsize \code{[K, 2]}, $K\le24$\\\scriptsize only at decode / IoU time};
  \draw[ar] (a)--(b) node[midway,above,lbl]{parser};
  \draw[ar] (b)--(c) node[midway,above,lbl]{decode};
\end{tikzpicture}
\end{center}

\begin{center}
\renewcommand{\arraystretch}{1.25}
\begin{tabular}{l l >{\raggedright\arraybackslash}p{5.4cm}}
\toprule
\textbf{Stage} & \textbf{Shape} & \textbf{Where} \\
\midrule
TFDS input        & \code{[N, 3972]} flat $xy$, $-1$ pad & dataset decoders \\
Training target   & \code{[N, 72]} radial & loss / parser \\
Prediction output & \code{[B, max\_det, 24, 3]} activated & detection generator \\
Cartesian (transient) & \code{[K, 2]} pixels & only at IoU/eval/viz \\
\bottomrule
\end{tabular}
\end{center}

\begin{pitfall}
In the model's \emph{prediction} output, the polygon \code{conf} values are already
sigmoid-activated by the detection generator --- they are \textbf{not} raw logits. Apply
your $0.4$ threshold directly. (And remember validity is keyed off the $-1.0$ sentinel,
not ``$\geq 0$'': a vertex can legitimately be slightly negative after a mosaic overflow
and is still real.)
\end{pitfall}

\section{Decoding a predicted polygon}
At inference, for each bin $i$ whose \code{conf}$_i \geq 0.4$:
\[
\text{angle} = (i + \text{offset}_i)\cdot 15^\circ, \qquad
\text{radius} = \text{dist}_i, \qquad
(x,y) = \text{box center} + \text{radius}\cdot(\cos,\sin)\,\text{angle}.
\]
Bins below threshold are dropped. The result is up to 24 Cartesian vertices, used for
visualization and for the polygon IoU metric.

\begin{teacher}
\textbf{Known limits of the format.} Because the origin is the box center and each bin
keeps only the \emph{farthest} crossing, the radial format cannot represent holes,
interior concavities, or shapes whose center lies outside the polygon. These are
properties of the representation itself, not bugs --- it trades exactness for a compact,
easy-to-predict encoding.
\end{teacher}
